\documentclass[11pt,a4paper]{book}

\usepackage[utf8]{inputenc}
\usepackage[T1]{fontenc}
\usepackage[top=2.5cm, bottom=1.5cm, left=1.5cm, right=1.5cm]{geometry}
\usepackage{amsmath,amssymb,amsfonts}
\allowdisplaybreaks
\usepackage{graphicx}
\usepackage{xcolor}
\usepackage{hyperref}
\usepackage{enumitem}
\usepackage[most]{tcolorbox}
\usepackage{fancyhdr}
\usepackage{array}
\usepackage{booktabs}
\usepackage{longtable}
\usepackage{multirow}
\usepackage{tikz}
\usepackage{float}
\usepackage{calligra}
\usepackage[export]{adjustbox}
\usepackage[font=small, labelfont=bf]{caption}
\usepackage{titlesec}

\usetikzlibrary{calc}

\definecolor{primaryblue}{HTML}{6BA3C8}
\definecolor{secondaryblue}{HTML}{A8D8E8}
\definecolor{lightgray}{HTML}{F5F5F0}
\definecolor{darkgray}{HTML}{3D3D3D}
\definecolor{dangerred}{HTML}{E8B4B8}
\definecolor{navyblue}{HTML}{4A3C3A}
\definecolor{chaptercolor}{HTML}{2F6FA3}
\definecolor{sectioncolor}{HTML}{5C524F}
\definecolor{subsectioncolor}{HTML}{C46A75}
\definecolor{subsubsectioncolor}{HTML}{7C68A6}
\definecolor{softlavender}{HTML}{D4C5E8}
\definecolor{softpeach}{HTML}{F4C2A3}
\definecolor{mintgreen}{HTML}{A8D5BA}
\definecolor{softcoral}{HTML}{F5A5A8}
\definecolor{brightred}{HTML}{FF2B2B}
\setlength{\headheight}{24pt}
\renewcommand{\arraystretch}{1.4}
\setlength{\tabcolsep}{6pt}
\newcolumntype{C}[1]{w{c}{#1}}
\newcolumntype{L}[1]{w{l}{#1}}

\let\originaltextbf\textbf
\renewcommand{\textbf}[1]{\textcolor{brightred}{\originaltextbf{#1}}}

\newenvironment{notelongtable}[1]{%
    \begingroup
    \setlength{\LTleft}{0pt}
    \setlength{\LTright}{0pt}
    \begin{longtable}{#1}
}{%
    \end{longtable}
    \endgroup
}

\renewcommand{\chaptermark}[1]{\markboth{Chapter \thechapter: #1}{}}
\renewcommand{\sectionmark}[1]{\markright{Section \thesection: #1}}

\pagestyle{fancy}
\fancyhf{}
\fancyhead[L]{%
    \begin{minipage}[t]{0.48\textwidth}
    \raggedright
    \adjustbox{max width=\linewidth}{\textcolor{darkgray}{{\bfseries\leftmark}}}%
    \end{minipage}}
\fancyhead[R]{%
    \begin{minipage}[t]{0.48\textwidth}
    \raggedleft
    \adjustbox{max width=\linewidth}{\textcolor{darkgray}{{\bfseries\rightmark}}}%
    \end{minipage}}
\fancyfoot[C]{\textcolor{darkgray}{\thepage}}
\fancyfoot[R]{\tiny\textcalligra{\textcolor{dangerred}{Miad}}}
\renewcommand{\headrulewidth}{2pt}
\renewcommand{\headrule}{\hbox to\headwidth{\color{primaryblue}\leaders\hrule height \headrulewidth\hfill}}

\fancypagestyle{plain}{%
    \fancyhf{}
    \fancyfoot[C]{\textcolor{darkgray}{\thepage}}
    \fancyfoot[R]{\tiny\textcalligra{\textcolor{dangerred}{Miad}}}
    \renewcommand{\headrulewidth}{0pt}
}

\setcounter{secnumdepth}{5}

\titleformat{\chapter}
  {\normalfont\Huge\bfseries\color{chaptercolor}}
  {Chapter \thechapter.}
  {1em}
  {}
  [\vspace{2pt}\textcolor{chaptercolor}{\titlerule[2.5pt]}]

\titleformat{\section}
  {\normalfont\huge\bfseries\color{sectioncolor}}
  {\thesection}
  {1em}
  {}

\titleformat{\subsection}
  {\normalfont\LARGE\bfseries\color{subsectioncolor}}
  {\thesubsection}
  {1em}
  {}

\titleformat{\subsubsection}
  {\normalfont\Large\bfseries\color{subsubsectioncolor}}
  {}
  {0pt}
  {\thesubsubsection\quad}

\titleformat{name=\subsubsection,numberless}
  {\normalfont\Large\bfseries\color{subsubsectioncolor}}
  {}
  {0pt}
  {}

\titleformat{\paragraph}[block]
  {\normalfont\large\bfseries\color{sectioncolor}}
  {}
  {0pt}
  {\theparagraph\quad}

\titleformat{name=\paragraph,numberless}[block]
  {\normalfont\large\bfseries\color{sectioncolor}}
  {}
  {0pt}
  {}

\titleformat{\subparagraph}[block]
  {\normalfont\normalsize\bfseries\color{darkgray}}
  {}
  {0pt}
  {\thesubparagraph\quad}

\titleformat{name=\subparagraph,numberless}[block]
  {\normalfont\normalsize\bfseries\color{darkgray}}
  {}
  {0pt}
  {}

\titlespacing*{\chapter}{0pt}{0pt}{2.3ex plus .2ex}
\titlespacing*{\section}{0pt}{3.5ex plus 1ex minus .2ex}{2.3ex plus .2ex}
\titlespacing*{\subsection}{0pt}{3.25ex plus 1ex minus .2ex}{1.5ex plus .2ex}
\titlespacing*{\subsubsection}{0pt}{3.25ex plus 1ex minus .2ex}{1.5ex plus .2ex}
\titlespacing*{\paragraph}{0pt}{2.5ex plus .5ex minus .2ex}{0.8ex}
\titlespacing*{\subparagraph}{0pt}{1.5ex plus .3ex minus .2ex}{0.5ex}

\newenvironment{definitionbox}[1][Definition]{%
    \begin{tcolorbox}[
        enhanced,
        breakable,
        colback=primaryblue!10,
        colframe=primaryblue!150,
        title=#1,
        fonttitle=\bfseries,
        sharp corners,
        boxrule=0.8pt,
        left=8pt,
        right=8pt,
        top=6pt,
        bottom=6pt
    ]
}{\end{tcolorbox}}

\newenvironment{formulabox}[1][Formula]{%
    \begin{tcolorbox}[
        enhanced,
        breakable,
        colback=softcoral!10,
        colframe=softcoral!200,
        title=#1,
        fonttitle=\bfseries,
        coltitle=white,
        sharp corners,
        boxrule=1pt,
        left=8pt,
        right=8pt,
        top=6pt,
        bottom=6pt
    ]
}{\end{tcolorbox}}

\newenvironment{infobox}[1][Information]{%
    \begin{tcolorbox}[
        enhanced,
        breakable,
        colback=mintgreen!10,
        colframe=mintgreen!200,
        title=#1,
        fonttitle=\bfseries,
        sharp corners,
        boxrule=0.8pt,
        left=8pt,
        right=8pt,
        top=6pt,
        bottom=6pt
    ]
}{\end{tcolorbox}}

\newcounter{qna}[chapter]
\renewcommand{\theqna}{\thechapter.\arabic{qna}}
\newenvironment{questionanswerbox}[1][Example]{%
    \refstepcounter{qna}
    \begin{tcolorbox}[
        enhanced,
        breakable,
        colback=softlavender!10,
        colframe=softlavender!200,
        title={#1~\theqna},
        fonttitle=\bfseries,
        coltitle=white,
        colbacktitle=softlavender!200,
        sharp corners,
        boxrule=1pt,
        left=8pt,
        right=8pt,
        top=6pt,
        bottom=6pt
    ]
}{\end{tcolorbox}}

\setlist[itemize,1]{label=\textcolor{primaryblue}{$\triangleright$}, leftmargin=*, itemsep=4pt, topsep=4pt}
\setlist[enumerate,1]{label=\textcolor{primaryblue}{\arabic*.}, leftmargin=*, itemsep=5pt, topsep=4pt}

\hypersetup{
    colorlinks=true,
    linkcolor=primaryblue,
    urlcolor=primaryblue,
    citecolor=primaryblue
}

\begin{document}
\setcounter{page}{1}
\tableofcontents
\clearpage

\chapter{Introduction}

% Source PDF page 5
\section{Information}

\begin{definitionbox}[Information]
Information refers to data or knowledge that is meaningful and has value in making decisions or solving problems.
\end{definitionbox}

\subsection{Importance of Information Theory}

\subsubsection{Why Information Theory Matters}
\begin{itemize}
    \item It helps us \textbf{quantify} information.
    \item It helps us \textbf{transmit} information efficiently.
    \item It helps us \textbf{optimize the delivery} of information in communication, computing, and biological systems.
\end{itemize}

\subsubsection{Real-Life Relevance}
\begin{itemize}
    \item Modern systems such as the \textbf{internet}, \textbf{mobile networks}, and \textbf{social media} rely on efficient data compression and error correction.
    \item These ideas ensure messages are transmitted \textbf{accurately} and \textbf{efficiently}.
    \item Information theory also helps reduce \textbf{noise interference} in signals.
    \item This improves reliability in applications such as \textbf{streaming}, \textbf{medical diagnostics}, and \textbf{genomic analysis}.
\end{itemize}

\subsection{Importance in Communication}

\subsubsection{Key Points}
\begin{itemize}
    \item Information theory provides mathematical frameworks to optimize \textbf{data transmission} and \textbf{storage}.
    \item It helps determine the maximum amount of information a channel can handle, that is, \textbf{channel capacity}.
    \item It helps determine how to encode messages so that transmission errors can be reduced.
\end{itemize}

\subsubsection{Communication Insight}
\begin{itemize}
    \item Cellular networks and internet protocols use \textbf{error correction} and \textbf{data compression}.
    \item These make communication \textbf{efficient} and \textbf{reliable}.
    \item Even in noisy environments, users can still experience:
    \begin{itemize}
        \item clear calls,
        \item fast internet speeds,
        \item accurate data exchange.
    \end{itemize}
\end{itemize}

\noindent This is why information theory is fundamental to seamless global communication.

\clearpage

% Source PDF page 6
\section{Discrete Source}

\begin{definitionbox}[Discrete Source]
A discrete source is a system that generates a \textbf{finite set of symbols or values} from a \textbf{predefined set of alphabets}. These symbols are discrete and distinct, such as letters, numbers, or pixels.
\end{definitionbox}

\subsubsection{Examples of Discrete Sources}
\begin{itemize}
    \item \textbf{Text:} A keyboard generates letters, numbers and special characters.
    \item \textbf{Audio:} Digital audio files sample sound at specific intervals, producing discrete data points.
    \item \textbf{Image/Video:} Images are composed of pixels with distinct color values, while videos are sequences of discrete image frames.
\end{itemize}

\subsection{Source Types}
\subsubsection{Memoryless Source}

\begin{definitionbox}[Memoryless Source]
A discrete source is memoryless if the symbols it generates are \textbf{independent} of each other.
\end{definitionbox}

\begin{itemize}
    \item The current symbol has \textbf{no dependence} on past symbols.
    \item Example: flipping a fair coin repeatedly produces independent outcomes $(H \text{ or } T)$.
    \item Such sources simplify analysis because each symbol can be treated individually without considering historical context.
    \item Example:A random number generator producing digits $(0\text{--}9)$ where each number is equally likely and independent of the previous numbers.
\end{itemize}

\subsubsection{Memory Source}

\begin{definitionbox}[Memory Source]
A discrete source is a memory source if the generation of a symbol \textbf{depends} on one or more \textbf{previous symbols}.
\end{definitionbox}

\begin{itemize}
    \item Such sources have \textbf{statistical dependency} between symbols.
    \item This dependency can be modeled using \textbf{Markov processes} or \textbf{conditional probabilities}.
    \item They are common in \textbf{natural language and encoded data streams}, where rules and structure influence symbol generation.
    \item Example: In English text, the letter `q' is almost always followed by `u', showing a strong dependency between these symbols.
\end{itemize}
\clearpage

% Source PDF page 7
\section{Self Information}

\begin{definitionbox}[Self Information]
Self information measures the \textbf{amount of uncertainty or surprise} associated with the occurrence of a specific event. It quantifies how much \textbf{information is gained} when an event occurs.
\end{definitionbox}

\begin{formulabox}[Self Information Formula]
Consider a discrete random variable $X$ with possible outcomes $x_i$, $i = 1, 2, 3, \dots, n$. Self information of the event $X = x_i$ is defined as
\[
I(x_i) = \log \frac{1}{P(x_i)} = -\log P(x_i).
\]
\begin{itemize}
    \item $I(x_i)$ = self information of the event $x_i$
    \item $P(x_i)$ = probability of the event $x_i$
\end{itemize}
\end{formulabox}

\subsection{Properties of Self Information}
\begin{enumerate}
    \item If we are absolutely certain of the outcome of an event even before it occurs, there is no information gained.
    \[
    I(x_i) = 0, \quad \text{if } P_i = 1
    \]
    This means \textbf{certainty carries no message information}.
    \item The occurrence of an event either provides some or no information, but never brings about a loss of information.
    \[
    I(x_i) \geq 0, \quad \text{for } 0 \leq P(x_i) \leq 1
    \]
    That implies \textbf{self information cannot be negative}.
    \item The \textbf{less probable an event is, the more information we gain} when it occurs.
    \[
    I(x_1) > I(x_2) \quad \text{if } P_1 < P_2
    \]
    \item For statistically independent events,
    \[
    I(x_1, x_2) = I(x_1) + I(x_2).
    \]
    Equivalently,
    \[
    I_A + I_B = \log \frac{1}{P_A} + \log \frac{1}{P_B}
    = \log \frac{1}{P_A P_B}
    = \log \frac{1}{P(AB)}.
    \]
\end{enumerate}

\subsection{Physical significance / insight of self information}
Consider a set of events $X = \{x_1, x_2, x_3\}$ with probability
\[
P = \{0.5, 0.25, 0.25\}.
\]

\bfseries{Step 1:} Self information for the first event:
\[
I_{x_1} = \log_2 \frac{1}{0.5} = -\log_2 0.5 = 1 \text{ bit}
\]

\bfseries{Step 2:} For the other two events:
\[
I_{x_2} = I_{x_3} = \log_2 \frac{1}{0.25} = -\log_2 0.25 = 2 \text{ bit}
\]

\begin{infobox}[Observation]
\begin{itemize}
    \item Higher probability requires fewer bits.
    \item Lower probability requires more bits.
    \item Lower probability therefore means a higher degree of uncertainty.
\end{itemize}
\end{infobox}

\clearpage

% Source PDF page 8
\section{Entropy}

\begin{definitionbox}[Entropy]
In information theory, entropy is the measurement of the average amount of information, that is, the number of bits of uncertainty, in a set of possible outcomes per symbol.
\end{definitionbox}

\subsection{Interpretation}
\begin{itemize}
    \item It measures the \textbf{unpredictability} of a random variable.
    \item It represents the minimum number of bits required to encode the information generated by a source.
\end{itemize}

\begin{formulabox}[Entropy of an Event]
\[
H(x_i) = -P_i \log_2 P_i.
\]
\begin{itemize}
    \item $H(x_i)$ = entropy contribution of the event $x_i$
    \item $P_i$ = probability of the $i$th event
\end{itemize}
\end{formulabox}

\subsection{Why This Form Appears}
\begin{itemize}
    \item Self information of an event is $\left(-\log_2 P_i\right)$.
    \item Entropy is an \textbf{average} information measure.
    \item So self information is weighted by its probability $(P_i)$.
\end{itemize}

\begin{formulabox}[Entropy of a Random Variable]
\[
H(X) = -\sum_{i=1}^{n} P_i \log_2 P_i
\quad \text{bit/symbol.}
\]
\begin{itemize}
    \item $H(X)$ = average entropy of the random variable $X$
    \item $n$ = total number of possible outcomes
    \item $P_i$ = probability of the $i$th outcome
\end{itemize}
\end{formulabox}

\begin{questionanswerbox}[Example]
A source generates information with probabilities
\[
P_1 = 0.1,\quad P_2 = 0.2,\quad P_3 = 0.3,\quad P_4 = 0.4.
\]
Find the entropy of the system and what percentage of maximum possible information is being generated by this source.

\textbf{Step 1: Calculate entropy}
\[
H(X) = -\sum_{i=1}^{n} P_i \log_2 P_i
\]
\[
= -\Bigl(0.1\log_2(0.1) + 0.2\log_2(0.2) + 0.3\log_2(0.3) + 0.4\log_2(0.4)\Bigr)
\]
\[
= 1.8464 \approx 2 \text{ bits/symbol.}
\]

\textbf{Step 2: Calculate maximum entropy}
\[
H_{\max} = \log_2 4 = 2.
\]

\textbf{Step 3: Calculate efficiency}
\[
\eta = \frac{H}{H_{\max}} \times 100
= \frac{1.8464}{2}\times 100
= 92.32\%.
\]
\end{questionanswerbox}

\begin{questionanswerbox}[Question]
\textbf{Q:} Why is \texorpdfstring{$P_i$}{Pi} multiplied with self information in the entropy equation?

\textbf{A:}
\begin{itemize}
    \item Entropy is an \emph{average} measure of information.
    \item So the self information term $\left(-\log_2 P_i\right)$ must be weighted by its probability $P_i$.
    \item This keeps highly probable events contributing more frequently and less probable events contributing less frequently to the final average.
    \item It also keeps the entropy value bounded properly for a valid probability distribution.
\end{itemize}
\end{questionanswerbox}

\clearpage

% Source PDF page 9
\subsection{Significance of Entropy}

\begin{itemize}
    \item \textbf{Measures uncertainty:} quantifies unpredictability in a source, providing insight into the information content.
    \item \textbf{Data compression:} sets the theoretical limits for lossless compression, helping design efficient encoding systems.
    \item \textbf{Error detection/correction:} guides the design of coding techniques to handle noise in communication channels.
    \item \textbf{System efficiency:} helps evaluate and optimize the performance of communication and storage systems.
    \item \textbf{Applications in AI/ML:} used in decision trees and deep learning for feature selection and model optimization.
\end{itemize}

\subsection{Relation between Hartley, Nats and Bits}

\subsubsection{Units of Self Information}
\begin{itemize}
    \item Self information in Hartley:
    \[
    I = \log_{10}\frac{1}{P}
    \]
    \item Self information in nats:
    \[
    I = \log_e \frac{1}{P}
    \]
    \item Self information in bits:
    \[
    I = \log_2 \frac{1}{P}
    \]
\end{itemize}

\begin{formulabox}[Hartley to Bits]
\[
1 \text{ Hartley}
= \frac{I}{-\log_{10} P}
= \frac{-\log_2 P}{-\log_{10} P}
= \frac{\log 10}{\log 2}
= \log_2 10
= 3.3219 \text{ bits.}
\]
\end{formulabox}

\begin{formulabox}[Hartley to Nats]
\[
1 \text{ Hartley}
= \frac{I}{-\log_{10} P}
= \frac{-\log_e P}{-\log_{10} P}
= \frac{\log 10}{\log e}
= \log_e 10
= 2.302 \text{ nats.}
\]
\end{formulabox}

\begin{formulabox}[Nats to Bits]
\[
1 \text{ Nat}
= \frac{I}{-\log_e P}
= \frac{-\log_2 P}{-\log_e P}
= \frac{\log e}{\log 2}
= \log_2 e
= \frac{1}{\ln 2}
= 1.442 \text{ bits.}
\]
\end{formulabox}

\begin{questionanswerbox}[Example]
If there are $100$ equiprobable outcomes, calculate the corresponding values in nats, bits and Hartley.

Given $N = 100$,
\begin{align*}
\text{In bits, } H(X) &= \log_2(100) = 6.64 \\
\text{In nats, } H(X) &= \log_e(100) = 4.605 \\
\text{In Hartley, } H(X) &= \log_{10}(100) = 2
\end{align*}
\end{questionanswerbox}

\clearpage

% Source PDF page 10
\subsection{Equal Probability Leads to High Entropy}

Consider a BPSK output
\[
x =
\begin{cases}
1, & \text{probability } p \\
0, & \text{probability } 1-p
\end{cases}
\]

\[
H(x) = -\sum_{x=0}^{1} P(x)\log_2 P(x)
\]
\[
= -p\log_2 p - (1-p)\log_2(1-p).
\]
\begin{itemize}
    \item $p$ = probability of symbol $1$
    \item $1-p$ = probability of symbol $0$
    \item $H(x)$ = entropy of the binary source
\end{itemize}

For equal probability $p=0.5$ and $1-p=0.5$,
\[
H(x) = -0.5\log_2 0.5 - 0.5\log_2 0.5 = 1 \text{ bit/symbol.}
\]

\begin{figure}[H]
    \centering
    \begin{tikzpicture}[scale=0.9]
        \draw[->] (0,0) -- (4.6,0) node[right] {$P$};
        \draw[->] (0,0) -- (0,3.1) node[above] {$H(P)$};
        \draw[smooth,domain=0.05:0.95,samples=120,thick,color=chaptercolor]
            plot (\x*4,{(-\x*ln(\x)/ln(2) - (1-\x)*ln(1-\x)/ln(2))*2.8});
        \draw[dashed] (2,0) -- (2,2.8);
        \node[below] at (0,0) {$0$};
        \node[below] at (2,0) {$0.5$};
        \node[below] at (4,0) {$1$};
        \node[left] at (0,2.8) {$1$};
    \end{tikzpicture}
    \caption{Entropy of a binary source versus probability}
\end{figure}

From the figure, it is clear that for $p=0$ or $p=1$, entropy $H(x)$ is zero, but for $p=0.5$, entropy is maximum. Equal probability leads to the highest uncertainty, so it also leads to the highest entropy.

\begin{questionanswerbox}[Example]
Calculate the entropy of the QPSK transmitted signal below:

\[
00\;\;01\;\;10\;\;10\;\;11\;\;11\;\;01\;\;01\;\;10\;\;00\;\;10\;\;10\;\;11\;\;01\;\;10
\]

QPSK has four distinct $(00, 01, 10, 11)$ bit combinations. Thus, $2^n = 4 \Rightarrow n = 2$. So, we have to take two bits for each symbol.

\begin{notelongtable}{|C{2.8cm}|C{2.2cm}|}
    \hline
    \textbf{Symbol} & \textbf{Times} \\ \hline
    \endhead
    \endfoot
    00 & 2 \\ \hline
    01 & 4 \\ \hline
    10 & 6 \\ \hline
    11 & 3 \\ \hline
\end{notelongtable}

\[
P_{00} = \frac{2}{15}, \qquad
P_{01} = \frac{4}{15}, \qquad
P_{10} = \frac{6}{15}, \qquad
P_{11} = \frac{3}{15}
\]

\[
H(X) = -\left[
P_{00}\log_2 P_{00}
+ P_{01}\log_2 P_{01}
+ P_{10}\log_2 P_{10}
+ P_{11}\log_2 P_{11}
\right]
\]
\[
= -\left[
\frac{2}{15}\log_2\frac{2}{15}
+ \frac{4}{15}\log_2\frac{4}{15}
+ \frac{6}{15}\log_2\frac{6}{15}
+ \frac{3}{15}\log_2\frac{3}{15}
\right]
\]
\[
= 1.889 \approx 2 \text{ bits/symbol.}
\]

\textbf{Physical significance:} Thus, we need at least $2$ bits to transmit the signal.
\end{questionanswerbox}

\clearpage

% Source PDF page 11
\section{Joint Entropy}

\begin{definitionbox}[Joint Entropy]
The joint entropy $H(X,Y)$ of a pair of random variables $(X,Y)$ with joint distribution $P(x,y)$ is the average uncertainty associated with the pair of outcomes taken together.
\end{definitionbox}

\begin{formulabox}[Joint Entropy Formula]
\[
H(X,Y) = -\sum_{x \in X}\sum_{y \in Y} P(x,y)\log_2 P(x,y)
\]
\[
= -E\log_2 P(x,y)
\]
\begin{itemize}
    \item $H(X,Y)$ = joint entropy of the variables $X$ and $Y$
    \item $P(x,y)$ = joint probability of the pair $(x,y)$
    \item $E[\cdot]$ = expectation
\end{itemize}
\end{formulabox}

\begin{questionanswerbox}[Example]
Two variables are considered for a student:
\begin{itemize}
    \item $X$ = study hours
    \item $Y$ = attendance
\end{itemize}

Their joint probabilities are listed below. Show that the variables are dependent on each other.

\begin{notelongtable}{|C{2.8cm}|C{2.8cm}|C{2.8cm}|}
    \hline
    \textbf{X} & \textbf{Y} & $\mathbf{P(X,Y)}$ \\ \hline
    \endhead
    \endfoot
    Low  & Low  & 0.20 \\ \hline
    Low  & High & 0.10 \\ \hline
    High & Low  & 0.15 \\ \hline
    High & High & 0.55 \\ \hline
\end{notelongtable}

\textbf{Step 1: Joint entropy contribution of each case}
\begin{align*}
H_{LL}(X,Y) &= -0.2\log_2(0.2) = 0.46 \\
H_{LH}(X,Y) &= -0.1\log_2(0.1) = 0.33 \\
H_{HL}(X,Y) &= -0.15\log_2(0.15) = 0.41 \\
H_{HH}(X,Y) &= -0.55\log_2(0.55) = 0.4743
\end{align*}

\[
H(X,Y) = 0.46 + 0.33 + 0.41 + 0.4743 = 1.6743
\]

\textbf{Step 2: Marginal entropy for study hours}
\[
P_L = 0.2 + 0.1 = 0.3, \qquad P_H = 0.15 + 0.55 = 0.7
\]
\[
H(X) = -\left[0.3\log_2(0.3) + 0.7\log_2(0.7)\right] = 0.881
\]

\textbf{Step 3: Marginal entropy for attendance}
\[
P_L = 0.2 + 0.15 = 0.35, \qquad P_H = 0.1 + 0.55 = 0.65
\]
\[
H(Y) = -\left[0.35\log_2(0.35) + 0.65\log_2(0.65)\right] = 0.934
\]
\end{questionanswerbox}

\begin{infobox}[Conclusion]
\[
H(X) + H(Y) = 0.881 + 0.934 = 1.815
\]
Since $H(X) + H(Y) > H(X,Y)$, the variables $X$ and $Y$ are dependent. We can also say that the variables contain mutual information.
\end{infobox}

\clearpage

% Source PDF page 12
\section{Mutual Information}

\begin{definitionbox}[Mutual Information]
Mutual information measures the amount of information that one random variable $X$ provides about another random variable $Y$. It quantifies the reduction in uncertainty about one variable when the other is known.
\end{definitionbox}

\begin{formulabox}[Mutual Information Formula]
\[
I(X;Y) = H(X) + H(Y) - H(X,Y)
\]
\[
\text{or equivalently,} \qquad I(X;Y) = H(X) - H(X \mid Y)
\]
\begin{itemize}
    \item $I(X;Y)$ = mutual information between $X$ and $Y$
    \item $H(X)$, $H(Y)$ = marginal entropies
    \item $H(X,Y)$ = joint entropy
    \item $H(X \mid Y)$ = conditional entropy of $X$ given $Y$
\end{itemize}
\end{formulabox}

\begin{questionanswerbox}[Example]
Calculate the joint entropy for the following data and also find the mutual information.

\begin{notelongtable}{|C{3.4cm}|C{2.4cm}|C{2.4cm}|}
    \hline
    \textbf{Y $\backslash$ X} & $\mathbf{0}$ & $\mathbf{1}$ \\ \hline
    \endhead
    \endfoot
    $0$ & $\dfrac{1}{2}$ & $\dfrac{1}{4}$ \\ \hline
    $1$ & $\dfrac{1}{8}$ & $\dfrac{1}{8}$ \\ \hline
\end{notelongtable}

\textbf{Step 1: Joint entropy}
\begin{align*}
H(X,Y) &= -\left[\frac{1}{2}\log_2\frac{1}{2} + \frac{1}{4}\log_2\frac{1}{4}
+ \frac{1}{8}\log_2\frac{1}{8} + \frac{1}{8}\log_2\frac{1}{8}\right] \\
&= 1.75
\end{align*}

\textbf{Step 2: Entropy of $Y$}
\[
P(Y=0) = \frac{1}{2} + \frac{1}{4} = \frac{3}{4}, \qquad
P(Y=1) = \frac{1}{8} + \frac{1}{8} = \frac{1}{4}
\]
\[
H(Y) = \frac{3}{4}\log_2\frac{4}{3} + \frac{1}{4}\log_2 4 = 0.811
\]

\textbf{Step 3: Entropy of $X$}
\[
P(X=0) = \frac{1}{2} + \frac{1}{8} = \frac{5}{8}, \qquad
P(X=1) = \frac{1}{4} + \frac{1}{8} = \frac{3}{8}
\]
\[
H(X) = \frac{5}{8}\log_2\frac{8}{5} + \frac{3}{8}\log_2\frac{8}{3} = 0.954
\]

\textbf{Step 4: Conditional entropies}
\[
H(X \mid Y) = H(X,Y) - H(Y) = 1.75 - 0.811 = 0.939
\]
\[
H(Y \mid X) = H(X,Y) - H(X) = 1.75 - 0.954 = 0.796
\]

\textbf{Step 5: Mutual information}
\[
I(X;Y) = H(X) + H(Y) - H(X,Y) = 0.954 + 0.811 - 1.75 = 0.015
\]
\end{questionanswerbox}

\clearpage

% Source PDF page 13
\section{Conditional Entropy}

\begin{definitionbox}[Conditional Entropy]
Conditional entropy measures the amount of uncertainty remaining in a random variable when another related variable is already known. It tells us how much additional information is still needed to describe X when Y is already provided.
\end{definitionbox}

\subsection{Simple Intuition}
Consider the set $\{1,2,3,4,5,6\}$. If we only know whether the outcome is odd or even, some uncertainty still remains.
\begin{itemize}
    \item If the outcome is \textbf{odd}, the possible values are $\{1,3,5\}$.
    \item If the outcome is \textbf{even}, the possible values are $\{2,4,6\}$.
\end{itemize}
This remaining uncertainty is the idea captured by conditional entropy.

\begin{formulabox}[Conditional Entropy Formula]
\[
H(Y \mid X) = -\sum_{x \in X} P(x)\sum_{y \in Y} P(y \mid x)\log_2 P(y \mid x)
\]
\[
= -\sum_{x \in X}\sum_{y \in Y} P(x,y)\log_2 P(y \mid x)
\]
\begin{itemize}
    \item $H(Y \mid X)$ = uncertainty of $Y$ after $X$ is known
    \item $P(y \mid x)$ = conditional probability of $Y=y$ given $X=x$
    \item $P(x,y)$ = joint probability of $X$ and $Y$
\end{itemize}
\end{formulabox}

\subsection{Relation with Joint Entropy}
Starting from the joint entropy,
\begin{align*}
H(X,Y) &= -\sum_{x \in X}\sum_{y \in Y} P(x,y)\log_2 P(x,y) \\
&= -\sum_{x \in X}\sum_{y \in Y} P(x,y)\log_2\bigl(P(x)P(y \mid x)\bigr) \\
&= -\sum_{x \in X}\sum_{y \in Y} P(x,y)\log_2 P(x)
- \sum_{x \in X}\sum_{y \in Y} P(x,y)\log_2 P(y \mid x) \\
&= -\sum_{x \in X} P(x)\log_2 P(x)
- \sum_{x \in X}\sum_{y \in Y} P(x,y)\log_2 P(y \mid x)
\end{align*}

Hence,
\[
H(X,Y) = H(X) + H(Y \mid X)
\]

Similarly,
\[
H(X,Y) = H(Y) + H(X \mid Y)
\]

\clearpage

% Source PDF page 14
\section{Relative Entropy}

\begin{definitionbox}[Relative Entropy]
Relative entropy, also called Kullback-Leibler divergence, measures how different an estimated probability distribution is from the true distribution.
\end{definitionbox}

\begin{formulabox}[Kullback-Leibler Divergence]
\[
D(P \parallel Q) = \sum_{x \in X} P(x)\log_2 \frac{P(x)}{Q(x)}
\]
\begin{itemize}
    \item $P(x)$ = true probability mass function
    \item $Q(x)$ = estimated probability mass function
    \item $D(P \parallel Q)$ = relative entropy from $Q$ to $P$
\end{itemize}
\end{formulabox}

\begin{infobox}[Basic Points]
\begin{itemize}
    \item If $P(x) = Q(x)$, then $D(P \parallel Q) = 0$.
    \item In the note's illustration:
    \[
    D(P \parallel P) = 0, \qquad D(P \parallel A) = 0.25, \qquad D(P \parallel B) = 1.25
    \]
    \item Higher relative entropy means the estimated distribution deviates more from the actual distribution.
\end{itemize}
\end{infobox}

\begin{questionanswerbox}[Example]
Consider a scenario where the true distribution is
\[
P_{\text{dog}} = 1, \qquad P_{\text{cat}} = 0
\]

For model $A$,
\[
P'_{\text{dog}} = 0.1, \qquad P'_{\text{cat}} = 0.9
\]

For model $B$,
\[
P'_{\text{dog}} = 0.9, \qquad P'_{\text{cat}} = 0.1
\]

Find the relative entropy using KL divergence.

\textbf{For model $A$:}
\[
D(P \parallel P'_A)
= P_{\text{dog}}\log_2\frac{P_{\text{dog}}}{P'_{\text{dog}}}
+ P_{\text{cat}}\log_2\frac{P_{\text{cat}}}{P'_{\text{cat}}}
\]
\[
= 1\log_2\frac{1}{0.1} + 0\cdot\log_2\frac{0}{0.9}
= \log_2 10
= 3.32 \text{ unit}
\]

\textbf{For model $B$:}
\[
D(P \parallel P'_B)
= P_{\text{dog}}\log_2\frac{P_{\text{dog}}}{P'_{\text{dog}}}
+ P_{\text{cat}}\log_2\frac{P_{\text{cat}}}{P'_{\text{cat}}}
\]
\[
= 1\log_2\frac{1}{0.9} + 0\cdot\log_2\frac{0}{0.1}
= 0.15 \text{ unit}
\]

So model $B$ is better because it has lower KL divergence.
\end{questionanswerbox}

\clearpage

% Source PDF page 15
\subsection{Mutual Information as Relative Entropy}
For two dependent random variables $(X,Y)$, mutual information can be written as the relative entropy between the joint distribution $P(x,y)$ and the product of the marginals $P(x)P(y)$:

\begin{formulabox}[Mutual Information from Relative Entropy]
\[
I(X;Y) = \sum_{x,y} P(x,y)\log_2 \frac{P(x,y)}{P(x)P(y)}
= D\bigl(P(x,y)\,\|\,P(x)P(y)\bigr)
\]
\end{formulabox}

\begin{align*}
I(X;Y)
&= \sum_{x,y} P(x,y)\log_2 \frac{P(x)P(y \mid x)}{P(x)P(y)} \\
&= \sum_{x,y} P(x,y)\log_2 \frac{P(y \mid x)}{P(y)} \\
&= \sum_{x,y} P(x,y)\log_2 P(y \mid x) - \sum_{x,y} P(x,y)\log_2 P(y) \\
&= -H(Y \mid X) + H(Y)
\end{align*}

So,
\[
I(X;Y) = H(Y) - H(Y \mid X)
\]

Also,
\[
H(X,Y) = H(X) + H(Y \mid X)
\]

Therefore,
\[
H(Y \mid X) = H(X,Y) - H(X)
\]

Substituting this into the previous expression,
\[
I(X;Y) = H(Y) - H(X,Y) + H(X)
\]

Thus the mutual information formula becomes
\[
I(X;Y) = H(X) + H(Y) - H(X,Y)
\]

\begin{figure}[H]
    \centering
    \begin{tikzpicture}[scale=0.95]
        \draw[thick, color=chaptercolor] (-1.1,0) circle (1.5);
        \draw[thick, color=chaptercolor] (1.1,0) circle (1.5);
        \fill[softlavender!35] (0,0) ellipse (0.4 and 1.15);
        \node at (-1.85,0) {$H(X \mid Y)$};
        \node at (1.85,0) {$H(Y \mid X)$};
        \node at (0,0) {$I(X;Y)$};
        \node at (-1.1,-1.95) {$H(X)$};
        \node at (1.1,-1.95) {$H(Y)$};
    \end{tikzpicture}
    \caption{Mutual information as the overlap between $H(X)$ and $H(Y)$}
\end{figure}

\end{document}
